\scanpage{31}
relation $(1)$ avec $l=k'-k=j'-j$ donne
\[
c_{j'-j,j+i}\,\xi_k=c_{j'-j,j}\,\xi_{k'\deleted{\unsure{+l|+i}}}
\]
et comme les coeffts sont des unités de $\mathbb Z_p$, on en conclut que $\xi_k=0\Longleftrightarrow\xi_{k'}=0$, cqfd.

\uline{Dém du lemme} 1\quad Considérons le développem$^{\rm t}$ $p$-adique de $i$
\[
 i=i_0+i_1p+\cdots+i_rp^r \qquad 0\leq i_\alpha\leq p-1\quad(0\leq\alpha\leq r)
\]
où on prendra $r$ arbitrairement grand (on ne suppose donc \uline{pas} que $i_r\neq0$ !). Considérons de $=$ le développem$^{\rm t}$ de $j=k-i$
\[
 j=j_0+j_1p+\cdots+j_rp^r\qquad 0\leq j_\alpha\leq p-1\quad(0\leq\alpha\leq r)
\]
L'hyp. $c_{i,j}\not\equiv0\ (p)$ signifie que
\[
 i_\alpha+j_\alpha\leq p-1\qquad\forall\ 0\leq\alpha\leq r
\]
Soit \deleted{\unsure{$\delta_\alpha$}}
\[
 \delta_\alpha=(p-1)-(i_\alpha+j_\alpha)\qquad 0\leq\delta_\alpha\leq p-1
\]
i.e.
\[
 j'_\alpha\overset{\text{déf}}{=}j_\alpha+\delta_\alpha=(p-1)-i_\alpha
\]
donc
\[
 0\leq j'_\alpha\leq p-1,
\]
et considérons
\[
 j'=\sum_0^r j'_\alpha p^\alpha
\]
donc $j'\geq j\ (\geq N)$ et
\[
 j'-j=\delta\overset{\text{déf}}{=}\sum_0^r\delta_\alpha p^\alpha.
\]
% margin note written sideways
\margin{\uline{Présentation plus simple} Soit $r\geq0$ tel que $p^r>\deleted{\unsure{\geq j|j}}i+j$ soit $\deleted{\unsure{\delta=p^r-1-}}\,\delta=p^r-1-(i+j)$ d'où $i+j+\delta=p^r-1$, $i+(j+\delta)=''$. Alors on a $\left\{\begin{aligned}c_{i,j+\delta}&\not\equiv0\ (p)\\c_{i+j,\delta}&\not\equiv0\ (p)\end{aligned}\right.$ (en plus de $c_{i,j}\not\equiv0\ (p)$) donc $\xi_k=0\Longleftrightarrow\xi_{k'}=0$, où $k'=j'+i=j+\delta+i=p^r-1$.}
On a alors \inserted{(en plus de $c_{i,j}\not\equiv0$)}
\[
\begin{aligned}
 c_{i,j'}&\not\equiv0\ (p) &&\text{car }i_\alpha+j'_\alpha\underset{\forall\alpha}{(=p-1)}\leq p-1\\
 c_{j'-j,j}&\deleted{\unsure{\not\equiv0|=}}=c_{\delta,j}\not\equiv0\ (p) &&\text{car }\delta_\alpha+j_\alpha=j'_\alpha\leq p-1
\end{aligned}
\]
% archive sheet number: 30

